\section{Hydrodynamics}
\label{sec:hydro}

\subsection{Conserved variables and the conservation laws}
\label{sec:hydro:euler}

The conserved state of a cell is the triple
%
\begin{equation}
  u = \bigl(\rho,\; \rho v,\; \mathcal{E}\bigr)^{\!\top},
  \qquad
  \mathcal{E} = \tfrac{1}{2}\rho v^2 + \rho e,
  \label{eq:cons}
\end{equation}
%
with $\rho e$ the internal energy density.  The map between $u$ and the
primitive triple $W = (\rho, v, p)$ is \texttt{UW\_conversions.f90}: the
density and the momentum are trivial, and the pressure follows from the
caloric equation of state of Sec.~\ref{sec:hydro:eos},
$p = p(\rho, \rho e)$, evaluated with the composition of the cell the state
belongs to.  For a reconstructed face state that cell is the cell the state
was extrapolated from, because the composition ratios the equation of state
needs are ratios of number densities and are therefore independent of the
density.

The equations integrated are the Euler equations in spherical symmetry with
gravity and a radiative source,
%
\begin{equation}
  \begin{aligned}
    \frac{\partial\rho}{\partial t}
      &+ \frac{1}{r^2}\frac{\partial}{\partial r}\bigl(r^2\rho v\bigr) = 0,\\
    \frac{\partial(\rho v)}{\partial t}
      &+ \frac{1}{r^2}\frac{\partial}{\partial r}\bigl(r^2\rho v^2\bigr)
       + \frac{\partial p}{\partial r}
       = -\rho\,\frac{\partial\Phi}{\partial r} + S_\mu,\\
    \frac{\partial \mathcal{E}}{\partial t}
      &+ \frac{1}{r^2}\frac{\partial}{\partial r}\bigl(r^2 (\mathcal{E}+p)v\bigr)
       = -\rho v\,\frac{\partial\Phi}{\partial r} + Q + S_{\mathcal{E}}
         - \frac{1}{r^2}\frac{\partial}{\partial r}\bigl(r^2 q_d\bigr),
  \end{aligned}
  \label{eq:euler}
\end{equation}
%
where $Q = \mathcal{H} - \Lambda$ is the net radiative heating minus cooling
per unit volume (Sec.~\ref{sec:thermal}) and $S_\mu$, $S_{\mathcal{E}}$ are the viscous and conductive
sources of Sec.~\ref{sec:hydro:transport}, both off unless a run asks for
them.  $q_d = \sum_s h_sJ_s$ is the interdiffusion enthalpy flux of a mixture
whose elements or transported carriers move relative to the mass-averaged
velocity $v$ \citep[eqs.~11--13]{Cook2009}: it exists where
\texttt{He\_diffusion} moves helium against hydrogen, where a molecular
carrier is transported and where \texttt{Ionization transport} carries the
ionization stages, and is written out in
Sec.~\ref{sec:mol:interdiffusion}: the element part is
\texttt{interdiffusion\_enthalpy\_face\_flux}
(\texttt{binary\_element\_diffusion.f90}), the carrier and stage parts are
\texttt{molecular\_carrier\_enthalpy\_face\_flux} and
\texttt{ionization\_stage\_enthalpy\_face\_flux}
(\texttt{diffusive\_photochemistry.f90}), and their divergences enter the
stationary energy row (\texttt{steady\_residual.f90}), the marching update
after an accepted element or carrier step (\texttt{EXHALE\_main.f90}) and the
advected-profile energy solve (Sec.~\ref{sec:postprocessing:adv}).  The factor $4\pi$ that would appear in the face areas and the cell
volume cancels between them and is simply not carried.

\subsection{The discrete divergence}
\label{sec:hydro:divergence}

\texttt{RK\_rhs.f90} assembles the right-hand side of \eqref{eq:euler} as a
flux difference \texttt{dF} and a source \texttt{S}, so that one forward
Euler stage of the marching loop is $u \leftarrow u - \Delta t\,(dF - S)$ and
the stationary residual is $R = dF - S$ with the radiative source subtracted
from its third row.  With
%
\begin{equation}
  A_\pm = r_{j\pm1/2}^2,\qquad
  \Delta V_j = \frac{r_{j+1/2}^3 - r_{j-1/2}^3}{3},\qquad
  \Delta r_j = r_{j+1/2} - r_{j-1/2},
  \label{eq:geom}
\end{equation}
%
and $F_\pm$ the numerical flux at the two faces of cell $j$
(Sec.~\ref{sec:numerics:riemann}), the three rows are
%
\begin{align}
  dF_1(j) &= \frac{A_+ F_+^{\rho} - A_- F_-^{\rho}}{\Delta V_j},
  \label{eq:dF1}\\
  dF_2(j) &= \frac{A_+ F_+^{\rho v} - A_- F_-^{\rho v}}{\Delta V_j}
             \;+\; w_p\,\frac{p_R - p_L}{\Delta r_j},
  \label{eq:dF2}\\
  dF_3(j) &= \frac{A_+ F_+^{\mathcal{E}} - A_- F_-^{\mathcal{E}}
             \;+\; \delta F_3}{\Delta V_j},
  \label{eq:dF3}
\end{align}
%
with $F^{\rho} = \rho v$, $F^{\mathcal{E}} = v(\mathcal{E}+p)$, and $p_L$, $p_R$ the face
pressures the Riemann solver returns at the lower and upper face.  The weight
$w_p$ and the content of $F^{\rho v}$ depend on the reconstruction, because
the spherical pressure term is split differently in the two schemes
(Sec.~\ref{sec:hydro:pressure}).

The term $\delta F_3$ in \eqref{eq:dF3} is the gravitational work,
%
\begin{equation}
  \delta F_3 = A_+ F_+^{\rho}\bigl[\Phi_i(j) - \Phi_c(j)\bigr]
             - A_- F_-^{\rho}\bigl[\Phi_i(j-1) - \Phi_c(j)\bigr],
  \label{eq:dF3p}
\end{equation}
%
with $\Phi_i$ the potential at the cell interfaces and $\Phi_c$ at the cell
centers (\texttt{set\_gravity\_grid.f90}).  Expanding the bracket separates
\eqref{eq:dF3} into the energy flux of \citet{Caldiroli2021}, their Eq.~(14),
$F^{\mathcal{E}} + F^{\rho}\Phi$ taken at the interface potential, minus
$\Phi_c(j)$ times the mass-flux divergence, which is their energy source
$\Phi_j L^\rho_j$.  The two forms are algebraically the same operator; the
code carries it as one term so that the mass flux the work term rides on is
the very flux the continuity row differenced.

\subsection{The spherical pressure term and gravity}
\label{sec:hydro:pressure}

The pressure gradient in spherical symmetry is
%
\begin{equation}
  \frac{\partial p}{\partial r}
    = \frac{1}{r^2}\frac{\partial}{\partial r}\bigl(r^2 p\bigr)
      - \frac{2p}{r},
  \label{eq:sphp}
\end{equation}
%
and the two reconstructions use the two sides of this identity.  Under PLM,
\texttt{Phys\_flux} adds $p$ to the momentum flux, so the first term of
\eqref{eq:sphp} is inside $dF_2$ and the geometric correction is carried by
the source; $w_p = 0$.  Under WENO3 the pressure is left out of the flux and
the whole gradient enters as the single face difference
$(p_R-p_L)/\Delta r_j$; $w_p = 1$.  \texttt{Source.f90} accordingly returns
%
\begin{equation}
  S_1 = 0,\qquad
  S_2 = -\frac{\rho_L + \rho_R}{2}\,
         \frac{\Phi_i(j) - \Phi_i(j-1)}{\Delta r_j}
        \;+\; \underbrace{\frac{A_+ - A_-}{\Delta V_j}\,p_c}
              _{\text{PLM only}},\qquad
  S_3 = 0,
  \label{eq:source}
\end{equation}
%
where $\rho_L$ and $\rho_R$ are the two reconstructed face densities that
bound the cell and $p_c$ is the cell's own pressure from the caloric equation
of state.  Because $(A_+-A_-)/\Delta V_j \to 2/r$ as the cell narrows, the
PLM branch of \eqref{eq:source} is the second term of \eqref{eq:sphp}.

The three physical terms of the momentum equation, the ram divergence, the
spherical pressure gradient and the weight $\rho\,\partial\Phi/\partial r$,
are gathered out of these pieces by
\texttt{momentum\_row\_terms\_of\_cell} (\texttt{RK\_rhs.f90}) and stored.  They are
not $dF_2$ and $S_2$: under PLM each of those holds a part of order $2p/r$
that cancels against the other, so a magnitude built from them alone reads
$2p/r$ on a uniform pressure at rest, where the physical force is zero.  The
gathered terms are what the stationary residual scales the momentum row by
(Sec.~\ref{sec:steady:residual}).

\subsection{The gravitational potential}
\label{sec:hydro:gravity}

\texttt{grav\_field.f90} carries two potentials, selected by
\texttt{Domain mode}.  The default is the full Roche potential of a
synchronously rotating planet,
%
\begin{equation}
  \Phi(r) = -\frac{b_0}{r}
            - \frac{b_0\,\mathcal{M}}{\tilde a - r}
            - \frac{b_0(1+\mathcal{M})}{2\tilde a^{3}}
              \left(\frac{\tilde a\,\mathcal{M}}{1+\mathcal{M}} - r\right)^{\!2},
  \qquad \mathcal{M} = \frac{M_\star}{M_p},\quad
  \tilde a = \frac{a_{\rm orb}}{R_0},
  \label{eq:roche}
\end{equation}
%
in code units, which is Eq.~(2) of \citet{Caldiroli2021}:
the planetary term, the stellar tidal term and the centrifugal term about the
center of mass.  With \texttt{Domain mode: Spherical} the last two terms are
dropped and $\Phi = -b_0/r$, which is what allows the domain to be carried
past the L1 radius.  The analytic derivative $d\Phi/dr$ is provided beside it
and is used by the initial-condition sonic-point search; the flux assembly
itself never differentiates $\Phi$, it differences the tabulated $\Phi_i$ and
$\Phi_c$, which are evaluated once at startup at the true face and center
radii.

The default outer radius in Roche mode is the L1 (Hill) radius
$r_{\max} = (3\mathcal{M})^{-1/3}\tilde a$; \texttt{Outer radius} overrides
it.

\subsection{Equation of state}
\label{sec:hydro:eos}

The thermal equation of state is that of an ideal gas of all particles
including free electrons,
%
\begin{equation}
  p = (n_{\rm tot} + n_e)\,k_B T
  \qquad\Longleftrightarrow\qquad
  T = \frac{p}{n_{\rm tot} + n_e}\ \ \text{in code units},
  \label{eq:thermal_eos}
\end{equation}
%
which is \texttt{comp\_T\_from\_p} and \texttt{comp\_p\_from\_T} of
\texttt{composition.f90}.  The particle and electron counts are built in one
place, \texttt{get\_species\_densities} in the same file, from the density and
the species mass fractions: \texttt{calc\_ntot} adds one particle for every
species the run carries (H\,I, H\,II, He\,I, He\,II, He\,III, every metal ion
stage, and each of $\mathrm{H}_2$, $\mathrm{H}_2^+$, $\mathrm{H}_3^+$,
$\mathrm{HeH}^+$, OH, $\mathrm{H_2O}$, CO), and \texttt{calc\_ne} adds the
free electrons weighted by the net charge of each species.  The
$\mathrm{He}\,2^3S$ column is not added by either: it is an excited level of
He\,I and its particle and its nucleus are already counted there.  Trace
metals enter both sums only under the \texttt{eos\_metals} policy of
Sec.~\ref{sec:mol:metals}, on by default, so that the legacy H/He-only budget
can be reproduced exactly.

The code never forms a mean molecular weight for the equation of state.  The
mean molecular weight of \citet{Caldiroli2021}, their footnote 1, is a derived
diagnostic here; \eqref{eq:thermal_eos} with the explicit particle count is
what the hydrodynamics uses, so the temperature of a cell is a function of its
composition and never of a fitted composition ratio.

The \emph{caloric} half of the equation of state is \texttt{caloric\_eos.f90}.
For an atomic gas it is the monatomic relation
$\rho e = p/(\gamma_{\rm ad}-1)$ with $\gamma_{\rm ad} = 5/3$, exactly as in
ATES.  This is not valid in a molecular layer: at the 400 to 1300~K of the
$\mathrm{H}_2$ base the rotational ladder of $\mathrm{H}_2$ is populated (the
$J=1$ level lies at 170~K), so an $\mathrm{H}_2$ molecule stores about
$\tfrac{5}{2}k_BT$ and the vibrational ladder adds to that above $\sim800$~K.
Where the gas is 80\% $\mathrm{H}_2$ the heat capacity is then wrong by a
factor 1.67.  The mixture relation used instead is
%
\begin{equation}
  \rho e = \tfrac{3}{2}(n_{\rm tot}+n_e)k_BT + n(\mathrm{H_2})\,k_B\,u_{\rm rv}(T),
  \qquad
  \frac{C_V}{k_B} = \tfrac{3}{2}(n_{\rm tot}+n_e) + n(\mathrm{H_2})\,c_{\rm rv}(T),
  \label{eq:caloric}
\end{equation}
%
so that the exponent making $\rho e = p/(\gamma_{\rm eff}-1)$ exact is
$\gamma_{\rm eff} = 1 + (n_{\rm tot}+n_e)k_B/C_V$, a function of both the
composition and the temperature.  The mean rovibrational energy $u_{\rm rv}$
and the heat capacity $c_{\rm rv}$ are evaluated from the complete bound
rovibrational ladder of $\mathrm{H}_2\,X^1\Sigma_g^+$, 302 levels to
51966~K, that is to the 4.478~eV dissociation limit
\citep[Table 2 through VizieR]{Roueff2019}; the same partition function is the
internal partition function of the
$\mathrm{H}+\mathrm{H}\leftrightarrow\mathrm{H}_2$ equilibrium constant
\eqref{eq:mol:r12}, so the energy and the chemistry cannot disagree.
$\mathrm{H}_2^+$, $\mathrm{H}_3^+$, $\mathrm{HeH}^+$ and the oxygen carriers
are counted as monatomic, a stated approximation valid while they stay trace
species (their largest ratio to $n(\mathrm{H_2})$ in the molecular layer of
the reference configuration is $1.7\times10^{-7}$).  Dissociation is not a
heat capacity and is not in \eqref{eq:caloric}; it is a chemical source
term.  The branch between the
two forms is structural, not a limit: a run with no molecules and a cell with
$n(\mathrm{H_2})$ exactly zero both take the constant-$\gamma_{\rm ad}$
expression verbatim.  The key \texttt{Caloric EOS: monatomic} forces the
monatomic form everywhere.

The mapping $\rho e \to p$ in the molecular branch is the inverse of
\eqref{eq:caloric} and costs one safeguarded Newton iteration per cell; the
mapping $p \to \rho e$ is direct.  The sound speed used by the Riemann
solvers and by the time step is $c_s = \sqrt{\gamma_{\rm eff}p/\rho}$
evaluated with the cell's own composition.

Two base composition scalars follow from the same module.  The gas mass per
hydrogen nucleus is
$m_{\rm H\text{-}nuc} = m(\mathrm{H\,I}) + m(\mathrm{He\,I})\,y + \sum_e A_e X_e$
with $y = n_{\rm He}/n_{\rm H}$; the base particle count per H+He nucleus is
%
\begin{equation}
  \hat n_{\rm bc} = \frac{1 + y + \sum_e X_e}{1 + y}
                    - \frac{1}{2}\,\frac{x_{\mathrm{H_2}}}{1+y},
  \qquad
  x_{\mathrm{H_2}} = \frac{2 q_{\mathrm{H_2}}(1+y)}{1 + q_{\mathrm{H_2}}},
  \label{eq:ntotbc}
\end{equation}
%
where $x_{\mathrm{H_2}}$ is the fraction of base hydrogen nuclei bound into
$\mathrm{H}_2$
and the last term is present only with \texttt{Molecular base: True}.  The
mixing ratio $q_{\mathrm{H_2}}$ is the one a lower-atmosphere handoff states,
or otherwise the chemical-equilibrium fit \eqref{eq:mol:qh2} at
$(p_{\rm base}, T_0)$ (Sec.~\ref{sec:mol:handoff}).  The same
$x_{\mathrm{H_2}}$ is imposed on the inflowing ghost composition, so the equation of state
and the species state describe one gas.

\subsection{The well-balanced option}
\label{sec:hydro:wb}

\texttt{Well balanced: True} (default off) changes the coordinate the
pressure is carried in.  The local hydrostatic equilibrium of cell $j$ is the
constant-density one through its own $(\rho_j, p_j)$,
%
\begin{equation}
  p_{{\rm eq},j}(r) = p_j - \rho_j\bigl[\Phi(r) - \Phi(r_j)\bigr],
  \label{eq:wb_eq}
\end{equation}
%
(Eq.~16 of \citealt{Kaeppeli2016}), whose two face values are
$P_{\rm up}(j) = p_j - \rho_j[\Phi_i(j)-\Phi_c(j)]$ and
$P_{\rm dn}(j) = p_j + \rho_j[\Phi_c(j)-\Phi_i(j-1)]$.  The reconstruction
then works on the departure from \eqref{eq:wb_eq} rather than on the pressure
itself, the Riemann pressure jump is formed as
$\Delta p = \Delta p_{\rm eq} + q_R - q_L$ from quantities that are all of
the size of the departure, and the momentum row is assembled from the face
pressure measured against each side's own equilibrium.  The two identities
that make the substitution exact are
%
\begin{equation}
  \begin{aligned}
    A_+P_{\rm up} - A_-P_{\rm dn} - (A_+-A_-)p_j
      &= -\rho_j\bigl\{A_+[\Phi_i(j)-\Phi_c(j)]
         + A_-[\Phi_c(j)-\Phi_i(j-1)]\bigr\}, \\
    P_{\rm up} - P_{\rm dn} &= -\rho_j\bigl[\Phi_i(j)-\Phi_i(j-1)\bigr],
  \end{aligned}
  \label{eq:wb_identities}
\end{equation}
%
the first for the PLM assembly and the second for WENO3.  Their left-hand
sides are the equilibrium part of the pressure terms the row carries and
their right-hand sides are the discrete gravitational source, so the two
cancel in the algebra, \texttt{Source.f90} returns $S_2 = 0$ and neither side
is ever evaluated.  This is the momentum source of
\citet{Kaeppeli2014}, their Eq.~(2.26), in spherical geometry.  What cancels is
still the physics the row balances, so its magnitude,
%
\begin{equation}
  \mathcal{W}_j = \frac{\bigl|\rho_j\{A_+[\Phi_i(j)-\Phi_c(j)]
        + A_-[\Phi_c(j)-\Phi_i(j-1)]\}\bigr|}{\Delta V_j}
  \quad\text{(PLM)},\qquad
  \mathcal{W}_j = \frac{\bigl|\rho_j[\Phi_i(j)-\Phi_i(j-1)]\bigr|}{\Delta r_j}
  \quad\text{(WENO3)},
  \label{eq:wb_weight}
\end{equation}
%
is formed separately and is the term the momentum residual is read against
under this option.

Under this option the lower boundary shares the same equilibrium: the
interior pressure its outgoing characteristic reads is $P_{\rm dn}(1)$, and
the two ghost cells continue \eqref{eq:wb_eq} below the level through the
face state (Sec.~\ref{sec:bcic:base}), so a column at rest in the discrete
equilibrium and meeting the reservoir's pressure at the level carries no
flux through the base face either.

\subsection{Low-Mach dissipation of the contact mode}
\label{sec:hydro:lowmach}

The HLLC flux resolves the contact wave exactly, so the dissipation it
applies to the entropy family is proportional to the contact speed $S_*
\simeq v$ and vanishes as the flow stalls; the acoustic families keep theirs
at $|v|\pm c_s$.  A stagnant layer can then carry a $2\Delta r$ pattern in
$v$ and $T$ that nothing in the scheme removes.  \texttt{Low-Mach damping:}
(default off, \texttt{low\_mach\_dissipation.f90}) adds a gated
fourth-difference stress \citep{Jameson1981} to the numerical momentum flux,
%
\begin{equation}
  D_p^{\,j+1/2} = \epsilon_4\,g(M)\,\rho_f\,\lambda_f
    \bigl(v_{j+2} - 3v_{j+1} + 3v_j - v_{j-1}\bigr),
  \qquad
  D_{\mathcal{E}}^{\,j+1/2} = v_f\,D_p^{\,j+1/2},
  \label{eq:lowmach}
\end{equation}
%
with $\rho_f$, $v_f$ the face averages and
$\lambda_f = \tfrac12[(|v|+c_s)_j + (|v|+c_s)_{j+1}]$ the face spectral
radius.  The gate
$g = [\max(0,\,1 - M_f^2/M_{\rm th}^2)]^2$, written in $M^2$ so that it is
smooth where the velocity changes sign, is unity at rest and identically zero
above $M_{\rm th}$ (default $10^{-3}$).  The pair $(D_p, D_{\mathcal{E}})$ conserves total
energy exactly and converts kinetic energy into heat and never the reverse;
the mass flux is untouched.  The term is added inside \texttt{RK\_rhs.f90}, so
the marching right-hand side and the stationary residual solve the same
equation, which is what distinguishes it from the Shapiro filter of
Sec.~\ref{sec:numerics:split}.

\subsection{Viscosity and heat conduction}
\label{sec:hydro:transport}

\texttt{Viscosity: True} and \texttt{Conduction: True} (both default off,
\texttt{viscous\_conduction.f90}) add the Navier-Stokes level of the same
equations, the level the CETIMB code (the time-dependent upper-atmosphere model of Koskinen and collaborators, \citealt{Koskinen2013, Koskinen2022}) carries.  For
the Newtonian stress in spherical symmetry the only nonzero components are
$\tau_{rr} = \tfrac{4}{3}\mu(v' - v/r)$ and
$\tau_{\theta\theta} = \tau_{\varphi\varphi} = -\tfrac12\tau_{rr}$, so
exactly
%
\begin{equation}
  S_\mu = \frac{1}{r^2}\frac{\partial}{\partial r}\bigl(r^2\tau_{rr}\bigr)
          + \frac{\tau_{rr}}{r},
  \qquad
  q_\mu = \frac{4}{3}\mu\left(\frac{\partial v}{\partial r} - \frac{v}{r}\right)^{\!2},
  \qquad
  S_{\mathcal{E}} = v S_\mu + q_\mu
        + \frac{1}{r^2}\frac{\partial}{\partial r}
          \left(r^2\kappa\frac{\partial T}{\partial r}\right).
  \label{eq:viscous}
\end{equation}
%
These differ from the expressions printed as Eqs.~(B5) and (B6) of
\citet{Koskinen2022}: those carry a different gradient term
($\mu' v'$ against $\tfrac43\mu' v/r$), a geometric coefficient $16/3$ against
$8/3$, and a dissipation that is not sign definite, and the printed pair does
not satisfy $vS_\mu + q_\mu = \nabla\!\cdot\!(\tau\cdot v)$.  Physical
correctness being the criterion, \eqref{eq:viscous} is what is implemented and
the discrepancy is recorded in the source.

Neither Koskinen paper prints the transport coefficients.  The heat
conduction coefficient is the number-weighted sum of the Banks \& Kockarts
(1973) coefficients of the constituents, as \citet{Salz2015} (their
Eqs.~24--26, electron--ion and neutral hydrogen) and \citet{Sutton2015}
(their Eq.~5, with helium) write it,
%
\begin{equation}
  \kappa = \frac{1}{n}\Bigl[\,n_e\,1.2\times10^{-6}\,T^{5/2}
     + n_{\rm HI}\,379\,T^{0.69} + n_{\rm HeI}\,299\,T^{0.69}
     + n_{\rm H_2}\,k_{\rm H_2}(T)\Bigr]
  \ \mathrm{erg\,cm^{-1}\,s^{-1}\,K^{-1}},
  \label{eq:transport_coeff}
\end{equation}
%
$n$ the number density of all particles and $T$ in K; ions contribute through
$n$ only.  The hydrogen term is the atomic hydrogen conductivity of
\citet{Watson1981}, $4.45\times10^{4}(T/1000\,\mathrm{K})^{0.7}$, to within
2 per cent between 300 and 7000~K; at He/H $=9.7$ the helium term decides
$\kappa$ (the coefficients stand in the ratio $299/379 = 0.79$; the mixture
value is 0.81 of the hydrogen one on the LHS~1140~b He/H\,$=9.7$ state,
MEASURED, update log stage~3 section~41), and it
agrees with the tabulated helium conductivity of \citet{Incropera2007}
(Table~A.4) to 1 per cent at 300 and 1000~K.  None of these sources gives an
H$_2$ coefficient: $k_{\rm H_2}(T) = A\,T^{s}\,[1 + c\,E(\theta_v/T)]$,
$E(x) = x^2 e^x/(e^x-1)^2$, $A = 272.523$, $s = 0.735652$, $c = 0.370446$,
$\theta_v = 5987$~K (the $v=0\to1$ spacing of H$_2$ in the
\citet{Roueff2019} ladder), an Eucken form fitted to the tabulated
conductivity of \citet{Incropera2007}, Table~A.4, over 200--2000~K (largest
deviation 2.5 per cent, rms 1.05 per cent); it is valid there and an
extrapolation outside (below 200~K, where the rotational heat capacity
freezes out, it overestimates, by 20 per cent at 100~K).  The Einstein function is evaluated as
$x^2e^{-x}/(1-e^{-x})^2$, which cannot overflow at low $T$, by its Taylor
series $1-x^2/12+x^4/240-x^6/6048$ below $x = 0.05$, where $1-e^{-x}$ loses
digits, and as zero above $x = 700$ (\texttt{einstein\_heat\_capacity}; the
$e^{x}$ form returned 0 below 16.9~K and NaN at 8.55--8.59~K until
2026-10-01); the H$_2$ term is formed only where $n_{\rm H_2}>0$.  The
molecular ions enter through $n$ and $n_e$ only.  A continuation factor
$0\le s\le1$ on $\kappa$ (\texttt{EXHALE\_CONDUCTION\_SCALE},
\texttt{conduction\_scale}; 1 when unset, any other value stops the run at
startup) multiplies the conductivity in the stationary row and the marching
update alike; a state solved at $s<1$ is a step of a continuation and not a
prediction of the model.  The factor is part of the state identity: wherever
$s\ne1$ the \texttt{cond} token of the options field carries $s$ itself to
seventeen significant digits (\texttt{factor\_token\_text},
\texttt{load\_IC.f90}), so two factors are two states.  The viscosity keeps the hydrogen value tied
to the Watson conductivity by the Chapman--Enskog/Eucken relation for a
monatomic gas, Prandtl number $2/3$,
$\mu(T) = \frac{4}{15}\frac{m_{\rm H}}{k_B}\,4.45\times10^{4}(T/1000\,\mathrm{K})^{0.7}$.  In code units
$\mu_{\rm code} = \mu_{\rm cgs}/(n_0 m_{\rm H} v_0 R_0)$, the inverse Reynolds
number, and $\kappa_{\rm code} = \kappa_{\rm cgs}T_0/(n_0m_{\rm H}v_0^3R_0)$
(equal to $\tfrac{15}{4}\mu_{\rm code}$ only for the Watson value the
viscosity is tied to, not for the mixture conductivity).

Both operators are assembled once as tridiagonal coefficient triplets on the
same faces and volumes as \eqref{eq:geom} (\texttt{viscous\_momentum\_coeffs},
\texttt{thermal\_conduction\_coeffs}), and the same triplets serve as the
source entering the stationary residual (\texttt{viscous\_conduction\_sources},
called by \texttt{steady\_residual.f90}, so every stationary solver of
Sec.~\ref{sec:steady} sees the terms) and as the matrix of a Crank-Nicolson
update in the marching loop (\texttt{viscous\_conduction\_step}, an
operator-split stage of each step), so the marched fixed point and the zero
of the residual are the same state.  The face coefficient is the cell
coefficient interpolated linearly in $r$ to the face, the same on both sides
of a face, so the operator is in flux form and its column sum telescopes to
the two boundary fluxes.

The outer boundary carries zero diffusive flux, which is the condition an
extrapolated outflow ghost admits.  At the inner boundary the viscous stress
is taken against the velocity of the anchored ghost, held fixed during the
implicit step.  The heat flux through the base face is taken against a
prescribed bath, not against the temperature the advective ghost happens to
carry: the lower atmosphere below the level is a heat bath on the time scale
of a stationary solution, whichever way the contact is upwinded.  The bath
temperature is \texttt{conduction\_base\_level\_T}, the lower atmosphere's own
temperature at the ghost center $r(0)$ (\texttt{base\_reservoir\_temperature\_at},
\texttt{base\_boundary.f90}): where a lower-atmosphere profile is handed over
(Sec.~\ref{sec:mol:handoff}) it is the profile's temperature at that radius,
linear in $r$ between its levels, and otherwise the reservoir's hydrostatic
isentrope continued from the level (Sec.~\ref{sec:bcic:base}); a reservoir
that returns no positive temperature leaves the ghost's own value.  The
conductivity of the ghost is evaluated at the same bath temperature, so the
coefficient and the temperature it acts on belong to one boundary condition.
The two choices agree at the level itself; one ghost cell down they differed
by 3.5\,K on the LHS~1140\,b photochemical profiles (188.89 against
185.42\,K), which moved the conductive flux into the level by about 4.6 per
cent at a fixed cell-1 temperature (MEASURED, update log stage~3
section~70).

The conduction operator is therefore conservative up to one named boundary
flux, the heat exchanged with the lower atmosphere through the base face,
\begin{equation}
  q_{\rm base} = -\kappa_{1/2}\,\frac{T_1 - T_{\rm bath}}{r_1 - r_0}
  \quad\text{(positive into the lower atmosphere)},
  \label{eq:qbase}
\end{equation}
the base face's own term of the first row of the triplets, sign reversed.
\texttt{heat\_conducted\_into\_lower\_atmosphere} forms it on the state the
certification measures (Sec.~\ref{sec:steady:certification}), that is on the state
about to be written and not on a perturbed Newton iterate, and
\texttt{write\_setup\_report.f90} writes it to \texttt{EXHALE\_resolved.out}
whenever conduction is on: \texttt{base\_conduction\_state}
(\texttt{unmeasured} before a certification has measured it, then
\texttt{measured}), \texttt{base\_conductive\_flux\_cgs} in
erg\,cm$^{-2}$\,s$^{-1}$, and the two temperatures of
\eqref{eq:qbase}, \texttt{base\_conduction\_T\_bath\_K} and
\texttt{base\_conduction\_T\_cell1\_K}.  The energy coupling of the
elemental-flux closure with the lower atmosphere reads this flux
(Sec.~\ref{sec:mol:handoff}).  On the LHS~1140\,b states it carried 35 to 52
per cent of the wind's absorbed heat (MEASURED, update log stage~3
section~76).

The Crank-Nicolson stage returns a state only when both tridiagonal
eliminations completed, every solved value is finite and every solved
temperature is at or above the floor of the equation of state (the floor of
the implicit energy step, Sec.~\ref{sec:thermal:tsolve}).  Otherwise it fails
with a named status and reason and returns no state
(\texttt{conduction\_report}).  The marching loop then refuses the whole
attempted step, in both run modes, exactly as it refuses a failed energy
update: the checkpoint is restored and the step is retaken at half the
interval, up to the retry budget of the attempted-step controller (eight
halvings, \texttt{n\_step\_retry\_max}); a stage that still fails at
$\Delta t/2^8$ stops the run with exit status~2 and the retry history
(Sec.~\ref{sec:numerics:rk}).  A shorter step is the remedy for the failure the
large step produces: a mode of the conduction operator with decay rate
$\lambda/\Delta t$ is multiplied by $(1-\lambda/2)/(1+\lambda/2)$, which is
negative for $\lambda>2$, so over a step longer than twice the diffusion time
of a cell the update overshoots, and halving $\Delta t$ halves $\lambda$.  The
floor value is never returned, because returning it would add energy with no
source behind it.

\subsection{The two-dimensional approximation}
\label{sec:hydro:2d}

A one-dimensional radial model of an irradiated sphere has to state how the
stellar flux and the mass loss are shared over the planet.  The mandatory key
\texttt{2D approximate method} carries that statement, in the four forms of
Table~\ref{tab:2d}.  The dilution factor $\xi$ is applied to every stellar
beam at the point where the beam is used, and in one place only
(\texttt{dayside\_dilution} in \texttt{parameters.f90}): the XUV photon grid,
the stellar Ly-$\alpha$ beam and the FUV bands all read it.  The stated
fluxes therefore keep their meaning as the flux at the planet's orbit, and
the correction to the mass-loss rate is applied once, at the end of the run
\eqref{eq:pp:mdot}.

\begin{table}[htb]
\centering
\caption{The \texttt{2D approximate method} options, as
\texttt{input\_read.f90} normalizes them, \texttt{dayside\_dilution} applies
$\xi$ and \texttt{EXHALE\_main.f90} applies the $\dot M$ factor.}
\label{tab:2d}
\begin{tabular}{llll}
\toprule
Key word & Stored as & $\xi$ & $\dot M$ factor \\
\midrule
\texttt{Mdot/4} & \texttt{Mdot/4}          & $1$    & $1/4$ \\
\texttt{Rate/2} & \texttt{Rate/2 + Mdot/2} & $1/2$  & $1/2$ \\
\texttt{Rate/4} & \texttt{Rate/4 + Mdot}   & $1/4$  & $1$   \\
\texttt{alpha}  & \texttt{alpha}           & $1$    & $1$   \\
\bottomrule
\end{tabular}
\end{table}

The first row is the recipe of \citet{Salz2016} that the ATES paper adopts:
the undiluted substellar rates with the resulting mass-loss rate divided by
four.  The second is the tidally locked reading, dayside averaged rates with
the outflow leaving one hemisphere; the third the rapidly rotating one,
rates reduced by four with the outflow leaving the whole surface.  These are
the three recipes of \citet{Caldiroli2021}, their Sec.~2.4.

The fourth is different in kind.  With \texttt{alpha} the run states a
coefficient $\alpha$ and the attenuated field carries the further factor
$1/(1+\alpha\tau)$ inside the column integral \citep{Sekiya1980, Odert2020},
so that the cell-mean attenuation becomes
%
\begin{equation}
  \bigl\langle e^{-\tau}\bigr\rangle_{\rm cell}
  = \int_0^1 \frac{e^{-\tau_{\rm out} - s\,\delta\tau}}
                  {1 + \alpha(\tau_{\rm out} + s\,\delta\tau)}\,ds,
  \label{eq:alpha2d}
\end{equation}
%
which is the cell-mean attenuation \eqref{eq:rad:cellmean} with the further
factor inside the integral; it has no closed form and is integrated by
composite three-point Gauss-Legendre on segments of optical depth at most
$0.25$, at most 256 of them (\texttt{cell\_mean\_attenuation},
\texttt{utilities.f90}).  With $\alpha = 0$, the default and the only value for which the field
is a pure exponential, it reduces to \eqref{eq:rad:cellmean} itself.

The range of validity is stated at the routine: $\xi = 1/2$ is the
illuminated hemisphere seen by an \emph{optically thin} absorber.  Where a
band is shielded the true shell average is lower, the slant columns away from
the substellar point being longer; it is 0.20 to 0.33 for the Lyman-Werner
band over a molecular layer.  No key asks for that refinement, so a shielded
band is diluted by the optically thin factor.
